\section{保护逻辑与 DER 故障穿越的可靠性接口}\label{sec:rel-pf}

\begin{theorynote}
本章围绕平台\emph{既有}动态（IEEE~1547 块）给出已落地的 \textbf{L1 轨迹类可靠性接口}：把继电器轨迹、
时限/反时限协调、断路器机械与 DER 故障穿越（FRT）事件，条件化为互斥的\emph{后果类}并折算为
EENS/LOLE/LOLF。给定轨迹入口消费代表性电流、DER 电压/相角轨迹和类条件切负荷阶段；在线入口由 Mass-Matrix DAE
生成故障电流与 DER 轨迹并反馈支路跳闸。自动拓扑、重合序列、瞬时闭锁和同步窗分别由可执行函数承担。
\end{theorynote}

\begin{boundarynote}
\textbf{实现状态（诚实边界，依据设计文档 \S18.1）}：下表区分给定轨迹、在线动态和拓扑后果的保真度。
除非结果显式报告 \fld{protection\_frt\_reliability\_coupled=true}，动态 IEEE~1547 与短路能力是\emph{独立}
分析，\emph{不}贡献该结果的 EENS/LOLE/LOLF/SAIDI 或恢复阶段 DER 可用性。
\begin{center}\footnotesize
\begin{tabular}{@{}L{5.7cm}L{2.5cm}L{5.5cm}@{}}
\toprule
能力 & 状态 & 可靠性含义\\
\midrule
DER 电压/频率越限计时、跳闸、重连合格与功率爬坡 & 已实现（选项，动态）& L1 类生成复用 IEEE~1547 状态机；通用 FMEA/MC 不自动调用\\
GFL/GFM 动态、电流限幅与优先级、volt-var/freq-watt & 已实现（选项，动态）& 供暂态研究，未转为列联类概率或 EENS 阶段\\
短路换流器贡献（GFL 电流源、GFM 阻抗后电压源）& 短路分析已实现 & 静态故障快照，非继电/FRT 事件演化\\
拒动/误动/拒开合、保护区扩展、重合概率 & 可靠性事件抽象已实现 & 作为\emph{输入}被故障模式 FMEA 与三阶段消费；启动值与选择性\emph{未}导出\\
反时限动作积分、复位、清除延时链、主/后备/断路器失灵事件树 & \textbf{L1 已实现} & 由给定轨迹生成互斥类和协调诊断\\
在线距离/差动、重合序列、DER--保护双向反馈与 DAE & \textbf{已实现} & 给定轨迹或在线 DAE 生成互斥类并进入年度事件树\\
\bottomrule
\end{tabular}
\end{center}
\end{boundarynote}

\begin{boundarynote}
\textbf{控制器—物理元件与部分可观测边界}：换流器控制器只与其所属 VSC/IBR 的端量传感器和电流/模式执行口关联。
保护功能不按故障元件强制一对一，而由保护区、CT/PT/状态接点量测集和受控断路器/接触器集共同定义；主后备可以区域重叠并动作于不同开断设备。
系统级恢复只能消费经可用信息通道送达的遥测、位置接点、保护事件和电源资格证书，不能直接读取 DAE 的真实完整状态。不同真实状态可产生相同报告；必需观测缺失时，遥控动作或电源资格不得准入，不得以内部真值补全。平台的有限 POMDP 是独立能力；除非某结果显式调用并声明该模型，不得将其冒充为通用可靠性恢复入口的当前行为。
\end{boundarynote}

\subsection{混合保护--穿越状态}\label{sec:rel-pf-state}
\begin{definition}[事件时状态]
事件时刻的混合状态为
\begin{equation}
 \zeta(t)=\big(x^{dyn}(t),\,\sigma(t),\,\chi(t),\,\kappa(t),\,\Omega(t),\,s^{br}(t),\,s^{rec}(t),\,\vartheta(t)\big),
\end{equation}
分别为电气/控制器微分状态、拓扑与并/离网模式、DER 连接态 $\chi_r\in\{0,1\}$、复合 DER 模式
$\kappa_r=(\kappa_r^{ctrl},\kappa_r^{FRT},\kappa_r^{lim})$、保护动作量/计时器 $\Omega_z$、断路器状态
$s_z^{br}$、重合器状态 $s_z^{rec}$ 与激活定值组 $\vartheta_z$。
\end{definition}
以电压角 $\delta_i$ 记电气量、以 $\vartheta_z$ 记保护定值，避免同符号歧义。事件间遵循微分代数
\begin{equation}
 \dot x^{dyn}=F_{\sigma,\kappa}(x^{dyn},y,u,w),\qquad 0=G_{\sigma,\kappa}(x^{dyn},y,u,w),
\end{equation}
$y$ 为网络代数量（电压、支路电流、接口功率），$u$ 为控制输入、$w$ 为扰动；离散守卫在
$g_e(\zeta(t^-))=0$ 时触发 $\zeta(t^+)=\mathcal R_e(\zeta(t^-))$，产生跳闸、模式切换、断路器动作、
重合与重连。复合模式使 GFL/GFM、穿越/跳闸、限幅/未限幅\emph{不}互斥：一台 DER 可同时为 GFM 构网、
强制穿越且电流受限。

\subsection{DER 电压/频率穿越自动机}\label{sec:rel-pf-frt}
设 DER $r$ 端正序电压 $V_r=|V_{i(r)}^{+}|$、量测频率 $f_r$。对配置带 $j$，其有效谓词 $h_{r,j}\in\{0,1\}$
与\emph{逐带}连续越限计时器
\begin{equation}
 D_{r,j}(t)=\begin{cases}t-\sup\big(\{s\in[t_0,t]:h_{r,j}(s)=0\}\cup\{t_0\}\big),&h_{r,j}(t)=1,\\[2pt]0,&h_{r,j}(t)=0,\end{cases}
\end{equation}
带条件清除即复位；这比把所有异常时间累加为一个标量严格，后者会错误合并不同激励与不同带。配置跳闸守卫
\begin{equation}
 \mathcal G_r^{trip}(t)=\bigvee_{j\in\mathcal J_r^{trip}}\big[D_{r,j}(t)\ge T_{r,j}^{trip}\big],\qquad
 T_{r,j}^{trip}=\mathcal T_{r,j}(\text{类别},\text{接入规程},\text{整定},\text{固件证据}).
\end{equation}
\emph{类别本身不决定跳闸曲线}：$T_{r,j}^{trip}$ 须落在适用要求内但由实际整定给出。模式转移与连接态为
\begin{equation}
 \kappa_r(t^+)=\Psi_r\big(\kappa_r(t^-),V_r,f_r,\{D_{r,j}\}_j,\eta_r^{device}\big),\quad
 \chi_r(t^+)=\begin{cases}0,&\mathcal G_r^{trip}=1,\\1,&\mathcal G_r^{reconnect}=1,\\ \chi_r(t^-),&\text{否则}.\end{cases}
\end{equation}
重叠带谓词须在 $\Psi_r$ 中给出显式优先级，否则两个同时激活的带会规定不兼容动作。

\subsection{穿越、瞬时闭锁与跳闸是不同状态}\label{sec:rel-pf-output}
DER 出力须按动态模式条件化：
\begin{equation}
 (p_r,q_r)=\begin{cases}
 (p_r^{cmd},q_r^{cmd}),&\kappa_r^{FRT}=\text{正常},\\
 (p_r^{FRT},q_r^{FRT}),&\kappa_r^{FRT}=\text{穿越},\\
 (p_r^{MC},q_r^{MC}),&\kappa_r^{FRT}=\text{瞬时闭锁},\\
 (0,0),&\kappa_r^{FRT}\in\{\text{跳闸},\text{闭锁},\text{重连等待}\},\\
 a_r(t)(p_r^{avail},q_r^{avail}),&\kappa_r^{FRT}=\text{爬坡}.
 \end{cases}
\end{equation}
瞬时闭锁策略 $(p_r^{MC},q_r^{MC})$ 是设备相关的（可为零、可保无功支撑），\emph{不得}假定所有并网 DER
都支撑电压，也\emph{不得}把瞬时闭锁等同于物理断开。$\Psi_r$ 须给出瞬时闭锁的进入/退出守卫（含恢复驻留、
迟滞、延迟升级为跳闸），否则穿越$\leftrightarrow$瞬时闭锁会在带边界颤振而终端类不唯一。

\subsection{故障期逆变器电流与 GFL/GFM}\label{sec:rel-pf-current}
公共电流限幅圆 $i_{d,r}^2+i_{q,r}^2\le(I_r^{max})^2$。有功优先时
\begin{equation}
 i_{d,r}=\operatorname{sat}(i_{d,r}^{ref},I_r^{max}),\quad
 i_{q,r}=\operatorname{sat}\!\Big(i_{q,r}^{ref},\sqrt{(I_r^{max})^2-i_{d,r}^2}\Big),
\end{equation}
无功优先时对调 $d/q$，其中 $\operatorname{sat}(x,\bar x)=\operatorname{sgn}(x)\min(|x|,\bar x)$。平衡三相
RMS 下故障期功率受\emph{电压相关}圆约束
\begin{equation}
 p_r^2+q_r^2\le(\sqrt3\,V_{LL,r}I_r^{max})^2\ \xrightarrow{\text{标幺}}\ p_r^2+q_r^2\le V_r^2(I_r^{max})^2,
\end{equation}
\emph{不}等于稳态铭牌圆；混用相/线电压与单相/三相功率会引入 $\sqrt3$ 或 $3$ 的因子错误。GFL 为 PLL 同步的
受控\emph{电流}源；GFM 为滤波/虚拟阻抗后的受控\emph{电压}源\emph{直至限幅改变有效模式}：
\begin{equation}
 \kappa_r=\begin{cases}(\mathrm{GFM},\kappa_r^{FRT},\text{电压构网}),&|i_r^{ref}|\le I_r^{max},\\
 (\mathrm{GFM},\kappa_r^{FRT},\text{电流受限}),&|i_r^{ref}|>I_r^{max}.\end{cases}
\end{equation}
保护到达与孤岛支撑必须用\emph{限幅后}模式；标签 \texttt{GFM} 本身不认证无限电压支撑或稳定孤岛。

\subsection{保护动作量：定时限与反时限}\label{sec:rel-pf-relay}
定时限（连续启动语义）$\dot D_z=1$（$\rho_z=1$）或 $-D_z/T_z^{reset}$（$\rho_z=0$），并在 $D_z\ge T_z^{delay}$
出令。反时限过流采用通用曲线并\emph{对动作量积分}以应对故障期电流变化：
\begin{equation}
 T_z^{curve}(M)=\mathrm{TMS}_z\!\left[\frac{A_z}{M^{p_z}-1}+B_z\right]+T_z^{add},\quad
 \dot\Omega_z=\begin{cases}1/T_z^{curve}(M),&M>1,\\ -\Omega_z/T_z^{reset},&M\le1,\end{cases}
\end{equation}
$M=|\widetilde I_z|/I_z^{pickup}$，出令时刻 $T_z^{relay}=\inf\{t:\Omega_z(t)\ge1\}$。当 DER 限幅/跳闸/电压
支撑或拓扑改变在继电器动作前改变电流时，该积分是\emph{必要}的。电气清除\emph{不}在出令时刻发生，清除链为
\begin{equation}
 T_z^{clear}=T_z^{relay}+T_z^{channel}+T_z^{tripcoil}+T_z^{mechanical}+T_z^{arc},
\end{equation}
$T_z^{channel}$ 仅在方案确用通信通道（导引/传输跳闸）时非零且条件于通信服务态；纯本地硬接线取 $0$。
所有时间量须用单一时基（通常秒），曲线系数须遵循选定 IEC/IEEE/厂商族而不混用。

\subsection{主后备协调与断路器失灵}\label{sec:rel-pf-coord}
对故障 $k$，$T_k^{pri}=\min_{z\in\mathcal P_k^{pri}}T_z^{clear}$、$T_k^{bak}=\min_{z\in\mathcal P_k^{bak}}T_z^{clear}$
（永不有效动作者取 $+\infty$）。须把\emph{物理清除结果}与\emph{协调诊断}分离：主清除事件 $\mathcal E_k^{pri}$
与满意协调 $\mathcal S_k^{coord}=\mathcal E_k^{pri}\cap\{T_k^{bak}-T_k^{pri}\ge\Delta T_k^{coord}\}$——主设备可
清除故障而后备裕度不足，这是\emph{带协调告警属性的主清除}，非未分类样本。跳令后电流持续即断路器失灵
\begin{equation}
 \mathcal F_z^{BF}=\big\{u_z^{trip}(t_z^{cmd}){=}1,\,s_z^{br}\ne\text{开},\,\inf_{\tau\in[t_z^{cmd},t_z^{cmd}+T_z^{BF}]}|\widetilde I_z(\tau)|>I_z^{open}\big\},
\end{equation}
启动失灵后备并扩展清除区；机械拒开概率是\emph{另一}按需失效，若已并入标定的失灵事件概率则\emph{不}再独立
相乘。原始事件可重叠，故须由显式事件树按优先级赋类：$(1)$ 若终端仍带电则\emph{拒清};$(2)$ 否则若任一非必要
健全区被跳则\emph{误协调};$(3)$ 否则按主/后备清除赋类。

\subsection{保护--FRT 后果类与 EENS}\label{sec:rel-pf-consequence}
把接口类扩展为六元组
\begin{equation}
 c=(c_D,c_I,c_P,c_{FRT},c_{rec},c_R),\qquad
 \pi_c(k)=\Pr(C_D{=}c_D,\dots,C_R{=}c_R\mid K{=}k),
\end{equation}
联合表示检测、隔离、保护、穿越、重连与恢复结果。设故障前暴露 DER 集 $\mathcal R_k$ 与终端分类时刻
$T_k^{term}=\min(T_k^{clear},T_k^{backup,max})$。以\emph{跳闸优先}定义互斥 FRT 划分
\begin{align}
 \mathcal E_k^{FRT,trip}&=\{\exists r\in\mathcal R_k:\chi_r(T_k^{term})=0\},\\
 \mathcal E_k^{FRT,MC}&=\{\chi_r(T_k^{term}){=}1\,\forall r\}\cap\{\exists r:\kappa_r^{FRT}(T_k^{term})=MC\},\\
 \mathcal E_k^{FRT,ride}&=\{\chi_r(T_k^{term}){=}1\,\forall r\}\cap\{\kappa_r^{FRT}(T_k^{term})\ne MC\,\forall r\};
\end{align}
$\mathcal R_k=\varnothing$ 时置 $c_{FRT}=\text{不适用}$ 而非依赖空真。\textbf{关键警示}：一台小光伏跳闸与一台
关键 GFM 储能跳闸\emph{不能}仅因同属 $\mathcal E_k^{FRT,trip}$ 而赋同一后果——聚合标签仅在 $\mathcal R_k$ 内
DER 后果等价、或物理引擎同时接收器件级轨迹 $\{\chi_r,\kappa_r,a_r\}$ 时才充分。类条件物理后果与 FRT 条件
EENS 增量为
\begin{equation}
\begin{aligned}
 S_{k,c}(t)&=S\big(\Phi_{k,c_P,c_{FRT},c_I}(G_0),\mathcal A_{c_R},
   \chi(t),\overline p^{avail}(t)\big),\\
 \Delta\mathrm{EENS}^{PFRT}&=\sum_k\lambda_k\!\sum_c\pi_c(k)
   \!\int_{\mathcal T_k}\!S_{k,c}(t)\,\mathrm dt,
\end{aligned}
\end{equation}
即\emph{动态保护/FRT 条件} EENS 与\emph{静态 DER 恒可用}基线之差。代表性类：主清除$\times$穿越/瞬时闭锁/
跳闸、后备清除、误协调、拒清（失灵为致因属性）、暂态自清除、重合闭锁、孤岛保护无效。

\subsection{微网孤岛可行性、储能续航与同步}\label{sec:rel-pf-microgrid}
DER 侧后果的\emph{岛级}闭合由物理微网模型给出。微网岛 $q$ 可服务须有可用构网锚
\begin{equation}
 y_q^{island}\le\sum_{r\in\mathcal R_q^{GFM}}x_r^p\,\phi_{ctrl(r)}(x^c),
\end{equation}
二值锚可用积已隐含“至少一个锚可用”，MILP 实现须为每个积 $x_r^p\phi_{ctrl(r)}$ 引入辅助二值。有功/无功
充裕度为 $\sum_{r}p_r^{DER}+p_q^{st}+p_q^{sh}=P_q^d+P_q^{loss}$、$\sum_{r}q_r^{DER}+q_q^{sh}=Q_q^d+Q_q^{loss}$，
换流器能力用二阶锥 $(p_r^{DER})^2+(q_r^{DER})^2\le(S_r^{rated})^2$；LP/MILP 须声明其多边形是\emph{保守内}逼近
\begin{equation}
 p_r\cos\theta_j+q_r\sin\theta_j\le S_r^{rated}\cos(\pi/J),\quad j=1,\dots,J,
\end{equation}
还是\emph{乐观外}松弛。\emph{无}额定、能量、无功能力与控制通道可用性的构网标志\emph{不}足以证明孤岛可维持。
储能续航按时序 SOC
\begin{equation}
 E_{b,t+1}=E_{b,t}+\eta_b^{ch}p_{b,t}^{ch}\Delta t-\frac{p_{b,t}^{dis}}{\eta_b^{dis}}\Delta t,\quad
 \underline E_b\le E_{b,t}\le\overline E_b,
\end{equation}
并以 $u_{b,t}\in\{0,1\}$ 禁止同时充放。\textbf{诚实}：储能\emph{不得}在每个恢复阶段或每次重叠停运独立按满
功率计入——事件阶段能量记账已在三阶段实现，多事件时序需序贯仿真。可再生受天气上限
$0\le p_{r,t}\le x_{r,t}^p\bar p_{r,t}^{weather}$；通信/控制丢失应按\emph{实际回退}模式（定值保持 /
本地下垂 / 失效跳闸）取值，而非一律当强迫停运——三种回退的可靠性价值迥异。合环前须满足同步窗
$|\Delta V|\le\Delta V^{max}$、$|\Delta f|\le\Delta f^{max}$、$|\Delta\theta|\le\Delta\theta^{max}$；忽略同步的
纯拓扑恢复计划须标注为可达服务的\emph{上界}。物理$\to$信息回耦（Level~3）中信息节点后备能量按
$B_{u,t+1}=\min(\overline B_u,\max(0,B_{u,t}+P_u^{ch}\Delta t-P_u^{load}\Delta t))$ 更新，可用度
$x_{u,t}^c=x_{u,t}^{intr}\mathbf 1[y_{\pi(u),t}{=}1\lor B_{u,t}{>}0]$，剩余续航 $T_u^{remain}=B_{u,t}/P_u^{load}$；
若 $T_u^{remain}<\Delta t$ 则节点在区间内失效，须事件队列或分数时长记账，端点二值更新会误判功能可用度。

\subsection{保真度阶梯与所需有效性标志}\label{sec:rel-pf-fidelity}
建议分级集成：$L0$——静态事件抽象（\S\ref{sec:rel-fa-cyber}），保护/FRT 为\emph{输入}概率与时长；
$L1$——对代表拓扑的给定故障/保护/FRT 轨迹生成互斥类，按场景权重折算 EENS；$L2$（当前离散事件子集已实现）——
在给定量测/DER 轨迹上联立多种继电判据、主后备/断路器和显式信息状态并自动生成三窗口后果；在线 $L2$ 由
Mass-Matrix DAE 生成网络/FRT/限幅轨迹；$L3$ 由连续信息物理事件队列执行年度抽样。类生成器
$\mathcal M_{PFRT}:(k,\omega)\mapsto(C_P,C_{FRT},C_{rec},T^{clear},\Phi_{k,c},\dots,\mathcal V^{PFRT})$
必须随附有效性向量 $\mathcal V^{PFRT}$（保真度、收敛、参数溯源、启动/FRT/限幅/同步边界附近的代理误差）。
严格地，最小混杂动态状态是 $(q,x)$，网络代数量 $y$ 仅在 $g_q(x,y)=0$ 的 index-1 正则分支上作为唯一
lifting；在线轨迹是在 hybrid time domain 上由模式内 Mass-Matrix DAE flow、可测 guard、首次到达时刻、
reset 和事件后代数重初始化构成的混杂执行。代数 fiber 为空/不唯一、$\partial g/\partial y$ 奇异或重初始化
残差不合格时必须进入 unresolved execution。阶段元组只是其可靠性奖励函数。完整执行 $\rho$ 是潜在生成对象，
类生成器只能消费有限记录 $\mathfrak o=\mathsf O_k(\rho)$；该记录包含已声明的量测/状态报告、事件词、监督命令和
计量或模型认证的服务边界，缺失字段保留 unknown，不得从隐藏 DAE 坐标补齐。有限类应定义为这些有限记录函数
在执行空间中的原像，并以 unresolved 原子闭合；$\mathsf O_k$ 非单射时，不同真实执行应留在同一类内。单个状态没有集合补集；
同层 class 的补集是其余同层类的不交并，子类在父类内的相对补集是 sibling classes 的不交并。该商空间只
保证同层互斥完备和跨层唯一父类，不宣称强 Markov lumpability 或连续轨迹双模拟；任何线性奖励的无偏压缩
都必须保存类条件矩，不能以一条名义代表轨迹替代。
这里的 $(q,x)$ 仅在固定完整场景后是最小动态状态，跨场景不自动构成 Markov state；若需要 Markov kernel，
必须加入足够的外生过程状态。当前求解器审计非线性收敛、耦合 Jacobian 分解和事件后代数残差，但不输出
$\inf_t\sigma_{\min}(\partial g/\partial y)>0$ 的全轨迹证书，因此 index-1 正则性是已选分支上的局部分析假设。
在线年度耦合器 \srcpath{compare\_online\_protection\_cyber\_reliability} 对同一故障分别运行主、后备
Mass-Matrix DAE 发现过程，保留 CT/PT 后的实测继电轨迹与 DER 逐点轨迹，再构造
$\mathcal M_{PFRT}$ 的输入并调用三级事件树。设 DAE 输出主保护命令时刻 $t_p^{cmd}$，断路器链总延时为
$d_p^{br}$，则事件树使用的主清除时刻严格取
\begin{equation}
 t_p^{clear}=t_p^{cmd}+d_p^{br};
 \end{equation}
不允许由 HTTP 或 GUI 再提供一个互相矛盾的清除时间。耦合成功时
\fld{online\_network\_dae\_coupled=true}，任一 DAE 不收敛、动作未施加或轨迹维数不一致均明确失败，
不得回退到用户给定轨迹。
\begin{boundarynote}
除非导出指标\emph{确实}消费上述耦合事件类或轨迹，任何单独动态研究都\emph{不}使可靠性结果 FRT 感知。合法置
\fld{protection\_frt\_reliability\_coupled=true} 需可追溯的类生成器；相关标志（设计文档 \S18.16）还包括
\fld{der\_ride\_through\_modelled}、\fld{protection\_coordination\_modelled}、\fld{breaker\_failure\_modelled}
等。只有 \fld{trace\_classes\_validated=true} 且类概率归一时聚合器才接受场景，并始终声明
\fld{online\_dae\_coupled=false}；在线年度耦合接口成功时单独置真，通用给定轨迹聚合器不冒充在线结果。
\end{boundarynote}

\subsection{case33bw AC/DC 恢复入口状态能力审计}
\label{sec:rel-pf-case33-propagation}
恢复模型中的二值换流器口径只读取
$a_c^{bin}=\mathbf 1[c\text{ authored in service}]$；DAE 入口口径则在故障
清除后读取端口在役状态与保护状态，
$a_c^{dae}=\mathbf 1[s_c^{port}=1\land s_c^{trip}=0]$。二者的差异数为
\begin{equation}
 N_{\Delta a}=\sum_c\left|a_c^{bin}-a_c^{dae}\right|.
\end{equation}
该比较遵循经典两状态可靠性元件模型与动态保护后果的分层定义，而不把额定容量
直接当成故障后可传输容量。可复现实现在
\srcpath{tools/case33bw\_acdc\_propagation\_study.cpp}，固定使用修正
\texttt{case33bw\_acdc} 的无故障、AC 母线 18 故障、传统 DC 母线 2 故障
以及启用直流故障控制的 DC 母线 2 故障轨迹。对故障量 $z_f$ 与相同初值的
无故障量 $z_0$，传播判据为
\begin{equation}
 \Delta z=\max_{t\in[t_f,T]}|z_f(t)-z_0(t)|,
 \qquad \Delta V>10^{-4}\ {\rm pu}\quad\text{or}\quad
 \Delta P_c>10^{-3}\ {\rm MW}.
 \label{eq:case33-propagation-gate}
\end{equation}
式~\eqref{eq:case33-propagation-gate} 是本研究预先声明的数值分离门，不是保护
定值；电压与功率门在 10 MVA 基准上均为 $10^{-4}$ pu，并远高于 $10^{-9}$
代数平衡容差。2 ms 结果的初始联合潮流残差为 $5.362\times10^{-11}$，
事件后最大代数残差为 $4.62\times10^{-14}$；无故障 AC/DC 电压漂移在输出
精度内为零。

AC 故障时，静态二值模型仍给出 VSC2 可用，而 DAE 在 0.365577673 s 记录
IEEE~1547 欠压跳闸并在恢复入口给出不可用，故 $N_{\Delta a}=1$。这证明
“二值在役”与“动态可恢复入口”可以不同。传统 GFL 路径下，DC 故障只产生
$2.1082\times10^{-10}$ pu 的 AC 电压偏差和 $1.2604\times10^{-10}$ MW
的换流器交流功率偏差，证明原路径实际上是单向耦合。启用
分段直流欠压降额/闭锁模型后，相同 DC 故障的
两项偏差分别增至 $8.807411\times10^{-2}$ pu 和 $1.05$ MW；1/2/5 ms
闭锁时刻为 $0.238/0.238/0.240$ s，传播分类一致，时刻离散误差范围为
2 ms。独立 Julia oracle 不调用 C++ 控制律，重新由 CSV 计算分段包络、20 ms
定时器、传统路径恒等性和传播门，全部通过。因此准入结论是：平衡相量
可靠性恢复入口具备 AC$\leftrightarrow$DC 传播，但不具备 EMT 直流故障电流
双向传播。

\begin{center}\footnotesize
\begin{tabular}{@{}p{3.0cm}p{4.2cm}p{3.3cm}p{3.2cm}@{}}
\toprule
\textbf{近五年 IEEE 证据} & \textbf{DC 故障至 AC 侧机制/验证} &
\textbf{与本文一致处} & \textbf{本文未覆盖处}\\
\midrule
Kaler--Yazdani 2022~\cite{relpf:kaler2022} & dq 动态模型、仿真及 1.25 kW
实验；故障闭锁拓扑限制电网侧电流 & 直流故障改变交流有功/电流，闭锁型 VSC
可用平均值控制描述 & 实验硬件拓扑与开关细节\\
Bandaru 等 2024~\cite{relpf:bandaru2024} & PSCAD/EMTDC；DC 故障时有功降至
零且可保留无功电压支撑 & 本文保留有功中断及其交流电压后果 & MMC 桥臂、
子模块与专用 FRT 策略\\
Yu 等 2024~\cite{relpf:yu2024}；Lv 等 2025~\cite{relpf:lv2025} & 控制/保护
分阶段动作；饱和改变故障响应 & 采用可审计的降额后定时闭锁顺序 & 桥臂电压
重构、饱和及阀级模型\\
Pandey 等 2025~\cite{relpf:pandey2025} & RTDS/CHIL；传统 BIC 二极管可形成
不控整流馈流，DC 故障显著改变 AC 端电流 & 证明 DC$\rightarrow$AC 物理因果
必须显式存在 & 二极管续流正是本文不支持的 EMT 路径\\
\bottomrule
\end{tabular}
\end{center}
表中论文验证的是机制而非本算例数值 oracle；不同拓扑、额定值和故障阻抗下的
波形幅值不可直接作为 case33 误差真值。本文阈值 $0.95/0.75$ pu 与 20 ms
延时均为研究参数，不是由文献反演或标准指定。模型仍不覆盖反并联二极管馈流、
MMC 子模块/桥臂饱和、开关纹波、DC 断路器电弧、直流过流保护及厂家解锁逻辑。

\subsection{N-k DAE 恢复入口类与既有切负荷 MILP}
\label{sec:rel-pf-nk-entry-class}
研究驱动 \srcpath{tools/nk\_acdc\_reliability\_study.cpp} 不把恢复优化结果
混入 trajectory class。对固定同时 N-k 进入事件 $e$，Mass-Matrix DAE 轨迹
$\rho_{e,\omega}$ 向恢复模型投影为附加不可用 VSC 集
\begin{equation}
 h_e(\rho_{e,\omega})=\mathcal U^{dae}_{vsc}(e,\omega),\qquad
 \rho_{e,\omega}\sim_e\rho_{e,\omega'}
 \Longleftrightarrow h_e(\rho_{e,\omega})=h_e(\rho_{e,\omega'}).
 \label{eq:rel-nk-entry-equivalence}
\end{equation}
运行状态 $\omega$ 的 AC/DC 负荷尺度仍进入每次恢复 MILP。原模型中的
$s_i^P,s_i^Q\ge0$ 分别进入有功/无功节点平衡，并由 $\lambda_{shed}$ 惩罚；
因此 class 只复用式~\eqref{eq:rel-nk-entry-equivalence} 的 DAE 入口，不删除切负荷
变量，不删除逐状态 MILP，也不使用类平均切负荷量替代状态后果。

该轨迹不是抽象状态序列。连续状态由每台换流器配置的 GFL（PLL、功率滤波、
电流响应/限幅及可选直流电容）或 GFM（内部相角/频率、电压/无功控制、功率滤波、
电流限幅及可选直流电容）状态组成；离散状态为开关位置、FRT、闭锁与保护状态。
交流侧由设备注入与 $Y_{ac}V_{ac}$ 的复电流平衡构成，直流侧由设备注入与
$G_{dc}V_{dc}$ 的电流平衡构成，VSC 交流功率按效率和直流电容能量平衡写入直流侧。
只有显式配置直流电压反馈、欠压降额或闭锁时，直流扰动才反向改变 GFL 交流功率/
电流指令；不能仅凭直流电容方程宣称 DC$\rightarrow$AC 传播。研究算例使用
MassMatrixDae 后向欧拉，并在每步内对控制状态、交流电压实虚部与直流电压做联立
阻尼 Newton 求解；状态事件经回退/二分定位、离散复位和事件后代数重初始化。
系统级恢复 MILP 不进入该 DAE，仅在 $T_{DAE}$ 后执行一次。

若后果相关 VSC 按 $(v_1,\ldots,v_{n_v})$ 排列，并令
$b^{(\ell)}=(b_{v_1},\ldots,b_{v_\ell})$，则可从包含全部已解析运行状态的单一集合开始，
逐个加入 VSC 可用状态。第 $\ell$ 层只能保持或细分第 $\ell-1$ 层，最终层等价于
式~\eqref{eq:rel-nk-entry-equivalence} 的完整不可用 VSC 集划分。层数因而等于
$n_v$，而不是固定三层；VSC 排序只影响中间层，不影响最终类。对固定进入事件，非空类互斥且
完备，类数不超过 $2^{n_v-|\mathcal V\cap\Delta\Phi_e|}$。按各类重新排列
$f_ep_\omega E_{e,\omega}$ 的求和严格保持事件频率和 EENS；这一恒等式不构成邻近运行状态
类别预测正确性的证据，预测仍必须由留一验证和被留出状态上的恢复重求解审计。

预测验证对每个事件留出一个运行状态，在归一化
$(\alpha_{AC},\alpha_{DC})$ 平面内从其余状态选择最近邻，只继承其入口类，随后
在被留出状态上重新求解恢复：
\begin{equation}
 \widehat S(e,\omega)=R\!\left(e,
 h_e(\rho_{e,j^*(e,\omega)}),\omega\right).
 \label{eq:rel-nk-loo-restoration}
\end{equation}
由类条件均值重构总体 EENS 是全概率公式恒等式，不作为式~\eqref{eq:rel-nk-loo-restoration}
的预测证据。成本审计保留恢复项：
\begin{equation}
 T_{direct}=N(\bar t_{DAE}+\bar t_{MILP}),\qquad
 T_{class}=T_{library}+N(\bar t_{lookup}+\bar t_{MILP}).
 \label{eq:rel-nk-cost}
\end{equation}
case33 的 92 个 N-1/N-2/N-3 联合进入事件与 6 个状态形成 552 个条件，全部求解；
92 个入口类给出 $1/6$ 压缩比，留一入口类准确率 100\%，恢复 EENS 误差为 0。
静态/DAE 入口 EENS 为 0.0905562/0.260831 MWh/yr；100000 样本的分项时间投影为
5.76--5.99 倍端到端加速，净节省成立但预注册 20 倍门槛失败。旧约 150 倍口径漏计了
每个状态的恢复 MILP，不能使用。

高阶故障组合和入口类划分不能混为同一层压缩。若 $d_e$ 为同时进入停运向量，
$\vartheta_e$ 为修复阶段参数，则恢复接口的完整离散特征为
\begin{equation}
 z_{e,\omega}=(d_e,h_e(\rho_{e,\omega}),\vartheta_e).
 \label{eq:rel-nk-complete-interface}
\end{equation}
当前代码以 `(event_index, class_label)` 分组，正好保持
式~\eqref{eq:rel-nk-complete-interface}：相同不可用 VSC 集在不同初始事件下不会合并。
对固定事件，$h_e$ 的非空逆像构成唯一最粗无损划分；任何保持恢复入口的其他划分都只能进一步细分它。
若全部枚举，事件数仍按 $\sum_m\binom{n_u}{m}$ 增长，trajectory class 只减少后续样本中重复 DAE 的
入口状态数，不消除故障组合空间。由于限流、孤岛和延时保护可导致故障阶数与跳闸集合非单调，不能未经
验证把父事件类传给子事件。当前驱动对 authored N-1/N-2/N-3 目录完全枚举；更高阶只能由具有频率的
联合事件模型/采样器给出，或在未展开 EENS 上界
$B_R=\sum_{(e,\omega)\in R}f_ep_\omega\bar E_{e,\omega}$ 小于给定容差后带证书停止。
该停止上界是理论扩展，当前算例没有把它宣称为已验证数值能力。

分域树模式曾因统一 AC+DC+VSC 连通图误剪 AC tie 而产生 150 个假不可行条件；
修正后统一图预剪仅用于统一树，分域模式由各域森林等式排环。该问题不是缺少
负荷切除能力。回归测试固定 VSC--DC--VSC 旁路、两条同时 AC 故障和两条必需
AC tie。当前频率为逐阶显式联合进入测度；顺序重叠、共因校准、EMT、原生三相
DAE 与恢复后非线性 AC/DC 认证均明确不覆盖。

case123 的 175 个联合进入事件与 6 个状态形成 1050 个条件；当前仅 892 个条件
进入并完成恢复 MILP，158 个在 DAE 或事件时代数网络求解阶段失败，未解析频率为
0.0319160/yr。已解析子测度上的静态/DAE 入口 EENS 为
0.0578627/0.235781 MWh/yr；已覆盖邻域入口类准确率为 100\%，但仍有
0.000761066/yr 的留一状态无已解析邻居。故准确性不准入。两次正式复跑的 8.81--8.91 倍成本投影也只作
诊断：失败长尾计入类库，而直接侧均值只来自已解析条件。未解析质量非零时
\fld{timing\_projection\_admissible=false} 且
\fld{net\_computational\_saving\_demonstrated=false}。这些失败属于 DAE 数值/模型
闭环，而非恢复 MILP 缺少负荷切除。

\subsection{与实现的对应}\label{sec:rel-pf-impl}
\begin{center}\footnotesize
\begin{longtable}{@{}L{5.4cm}L{9.0cm}@{}}
\toprule
\textbf{理论对象} & \textbf{平台实现状态}\\
\midrule
\endfirsthead
\toprule
\textbf{理论对象} & \textbf{平台实现状态}\\
\midrule
\endhead
\bottomrule
\endfoot
静态保护事件抽象（拒动/误动/拒开合/区扩展）& \textbf{已实现}：故障模式 FMEA 与三阶段消费为输入概率与配置时长；\srcpath{src/reliability/failure\_mode.cpp:build\_consequence\_patch}\\
短路换流器贡献（GFL 电流源、GFM 阻抗后电压源）& \textbf{已实现}（短路模块，静态快照）；\srcpath{src/short\_circuit/}\\
DER IEEE~1547 越限/跳闸/重连状态机 & \textbf{已实现}（动态，选项），由 L1 轨迹分类器复用；\srcpath{src/dynamics/}\\
逐带 FRT 计时器 / 模式自动机 $\Psi_r$ / 瞬时闭锁 & \textbf{已实现}：IEEE~1547 逐带计时与 \srcpath{DERMomentaryCessationSettings}\\
定/反时限、方向闭锁、距离/差动、清除链 / 主后备协调 / 断路器失灵 & \textbf{给定轨迹 L2 已实现}；\srcpath{src/reliability/protection\_frt.cpp:evaluate\_protection\_relay / evaluate\_protection\_coordination}\\
共享信息路径、QoS、共因环境、供电与电池自治、功能绑定 & \textbf{离散状态 L2 已实现（最多 20 个二值信息元件）}；\srcpath{src/reliability/protection\_frt.cpp:enumerate\_protection\_cyber\_states}\\
静态 FMEA / 仅保护 / 信息物理联合指标对比 & \textbf{已实现}；\srcpath{src/reliability/protection\_frt.cpp:compare\_protection\_cyber\_reliability}\\
保护--FRT 后果类 $\pi_c(k)$ 与 EENS/LOLE/LOLF & \textbf{L1 已实现}（轨迹验证类 + 外部切负荷阶段）；\srcpath{src/reliability/protection\_frt.cpp:aggregate\_protection\_frt\_reliability}\\
在线 DAE、瞬时闭锁、自动后果拓扑与重合/同步反馈 & \textbf{已实现}：\srcpath{online\_protection.cpp}、\srcpath{simulate\_recloser\_fuse\_sectionalizer}、\srcpath{evaluate\_automatic\_protection\_topology}、\srcpath{evaluate\_microgrid\_synchronization}\\
\end{longtable}
\end{center}

\subsection{参考文献}
\begin{thebibliography}{99}
\bibitem{relpf:ieee1547} IEEE Std 1547-2018, \emph{IEEE Standard for Interconnection and Interoperability
 of Distributed Energy Resources with Associated Electric Power Systems Interfaces}, 2018.
\bibitem{relpf:iec60255} IEC 60255-151, \emph{Measuring relays and protection equipment -- Part 151:
 Functional requirements for over/under current protection}, 2009.
\bibitem{relpf:c37} IEEE Std C37.113, \emph{IEEE Guide for Protective Relay Applications to Transmission Lines}.
\bibitem{relpf:kundur} P. Kundur, \emph{Power System Stability and Control}. New York, NY, USA:
 McGraw--Hill, 1994.
\bibitem{relpf:yazdani} A. Yazdani and R. Iravani, \emph{Voltage-Sourced Converters in Power Systems}.
 Hoboken, NJ, USA: Wiley--IEEE Press, 2010.
\bibitem{relpf:kaler2022} S. Kaler and A. Yazdani, ``A DC-side fault-tolerant
 bidirectional AC--DC converter for applications in distribution systems,''
 \emph{IEEE Access}, vol. 10, pp. 52625--52637, 2022,
 doi: 10.1109/ACCESS.2022.3171346.
\bibitem{relpf:bandaru2024} S. Bandaru \emph{et al.}, ``AC and DC fault
 ride-through strategies for hybrid MMC based HVDC link and interaction with
 multimachine power system dynamics,'' \emph{IEEE Trans. Ind. Appl.}, 2024,
 doi: 10.1109/TIA.2024.3396790.
\bibitem{relpf:yu2024} X. Yu \emph{et al.}, ``An active DC fault current
 limiting control for half-bridge MMC based on arm voltage reconstruction,''
 \emph{IEEE Trans. Power Del.}, 2024, doi: 10.1109/TPWRD.2023.3294173.
\bibitem{relpf:lv2025} Z. Lv \emph{et al.}, ``Modeling of the DC fault response
 of MMCs with fault control capability considering converter saturation,''
 \emph{IEEE Trans. Power Del.}, 2025, doi: 10.1109/TPWRD.2024.3485910.
\bibitem{relpf:pandey2025} R. Pandey \emph{et al.}, ``Time-domain backup
 protection for converter dominated hybrid AC--DC microgrids,''
 \emph{IEEE Trans. Ind. Informat.}, 2025, doi: 10.1109/TII.2025.3574424.
\bibitem{relpf:designdoc} HySim 团队，\emph{可靠性数学模型与智能信息物理评估}，\S18（设计文档）。
\end{thebibliography}
